% sec-07-ai-peer-review.tex - Section 7, AI Peer Review.  A main section.
% DRAFT stage.  Four figure slots, Figures 21 to 24, and four tables, Tables 21
% to 24.  Primary source is the AI peer review archive under
% new-trial-system/inputs; the funding-application review schedule comes from
% funding/RFA-RM-27-001-v2.
\section{AI Peer Review}
\label{sec:peerreview}

\subsection{The study the method comes from}

The AI peer review method used throughout this repository was established in a
glioblastoma study deposited November 30, 2025 \cite{aiprgliomatc}. That study
was about clinical trial patient matching in glioblastoma, not pancreatic
cancer, and the disease is not the point: what transfers is the review method,
and it transfers without modification to the pancreatic cancer work described in
sections 3 through 6.

\draftinstr{Stage 7: quote directly from
\appfile{new-trial-system/inputs/AI\_Peer\_Review\_Acceleration\_of\_LLM\_
Generated\_Glioblastoma\_Clinical\_Trial\_Patient\_Matching\_ML\_\_FDA\_ICH\_
ISO\_\_and\_FastAPI.zip}, main.tex, the abstract sentence beginning ``A primary
benefit of AI peer review is the ability to appropriately conduct iterative
analyses and corrections throughout the entire manuscript process'' and the
sentence that follows it. State explicitly that this was a glioblastoma study
and that the method applies to the pancreatic cancer study.}

\figslot{mermaid-type / sequenceDiagram}{9.4}{Figure 21. One manuscript on two
clocks: the prior system's serial human\\rounds over months, and the new system's
parallel model rounds within hours.}

Figure 21 lays the two regimes on one clock. The prior system's best-case review
cycle is seven to eight weeks and its typical journal processing several months
across one or two rounds; faster online journals reach about one month. The new
system's round closes the same day, and it closes during development rather than
after completion.

\draftinstr{Stage 7: build Table 21 from
\appfile{new-trial-system/d2/fig-22-peer-review-economics-grid.md} cell value
table. Columns: Axis, Prior system, New system, Ratio. Seven rows, the last
being corrections applied before release.}

\begin{apptable}
\begin{tabularx}{\textwidth}{>{\raggedright\arraybackslash}p{3.9cm} >{\raggedright\arraybackslash}p{3.6cm} >{\raggedright\arraybackslash}p{3.4cm} >{\raggedright\arraybackslash}X}
\headrow \thead{Economic axis} & \thead{Prior system} & \thead{New system} & \thead{Ratio}\\
\midrule
Latency, one round & 7 to 8 weeks best case & Same day & About 50 to 1 \\
Typical journal processing & Several months, 1 to 2 rounds & Hours, rounds unlimited & About 300 to 1 \\
Reviewers per round & 2 to 3 humans & 3 model manufacturers & 1 to 1 by count \\
Entry point & After completion & During development & Not comparable \\
Marginal cost per round & Reviewer time, unpriced & Inference cost, tens of dollars & Orders of magnitude \\
Artifacts absorbed per year & 2 to 6 per author & Over 30 deposited in 2026 & About 6 to 1 \\
Corrections before release & None, review follows release & Every round, before deposit & The decisive axis \\
\end{tabularx}
\end{apptable}
\vspace{-0.6cm}
\tabcap{Table 21. Six economic axes of one review round under both regimes, and\\the
seventh axis on which the two are not comparable at all.}

\figslot{d2-type / grid}{6.4}{Figure 22. Six economic axes of one review round
under both regimes, with\\the ratio in each cell and the volume the prior system
would have to absorb.}

Figure 22 carries the same six axes as a grid, with the ratio computed in each
cell.

\subsection{Three manufacturers, one human gate}

\draftinstr{Stage 7: quote directly from
\appfile{funding/RFA-RM-27-001-v2/LaTeX Source Files.zip},
sections/ai-peer-review-context.tex, the paragraph assigning Claude Code the
primary production role and ChatGPT Codex and Google Gemini the two independent
reviewer roles, and the closing sentence stating that these systems are not
applicants, investigators, sponsors, regulators, clinicians, or decision-makers,
and that the PI remains responsible for source verification, regulatory
accuracy, scientific interpretation, participant protection, and final
submission.}

\begin{apptable}
\begin{tabularx}{\textwidth}{>{\raggedright\arraybackslash}p{2.6cm} >{\raggedright\arraybackslash}p{3.1cm} >{\raggedright\arraybackslash}X}
\headrow \thead{Manufacturer} & \thead{Role} & \thead{What that role covers in this work}\\
\midrule
Anthropic & Primary production & Code, regulatory adaptations, protocols, IND materials, diagrams, repository-scale artifacts \\
OpenAI & First independent reviewer & Verification, validation, uncertainty quantification, external standards, trial-literature synthesis, deep research, PDF review \\
Google & Second independent reviewer & Accelerated code problem solving, meta-verification of prior stages, administrative workflow support \\
Human PI & Final authority & Source verification, regulatory accuracy, scientific interpretation, participant protection, final submission \\
\end{tabularx}
\end{apptable}
\vspace{-0.6cm}
\tabcap{Table 22. The three model manufacturers and the human principal\\investigator,
and the role each holds in the reviewed pancreatic cancer work.}

Table 22 names the three manufacturers, the role each holds, and the human
principal investigator who retains final authority over all three.

\figslot{plantuml-type / activity with fork and join}{10.2}{Figure 23. One
artifact reviewed concurrently by three model manufacturers,\\joined into one
recorded disagreement set, and cleared only by the human PI.}

Figure 23 shows where the concurrency sits and where it stops. Three
manufacturers read one frozen artifact and its recorded hash, so no branch can
observe another's output and a disagreement is evidence rather than an artifact
of ordering. The join produces three reports over one object, and the human
principal investigator is the only gate past it.

\subsection{What the review actually found}

\draftinstr{Stage 7: report the glioblastoma study's measured outcomes from the
same archive: initial test-set performance 67.3 percent on a BioClinicalBERT
pipeline; a triple review by three manufacturers all recommending a tabular
model; the TabPFN-v2 replacement; test-set accuracy 94.0 percent after the
recommendations were implemented; an evaluation standard scoring the workflow at
87.5 percent overall quality and 0.831 efficiency; eleven unsuccessful FastAPI
attempts followed by a successful Opus and Sonnet strategy; and three regulatory
analyst scores of 8.0, 9.1 and 9.1 out of 10 with the reason each gave.}

\begin{apptable}
\begin{tabularx}{\textwidth}{>{\raggedright\arraybackslash}p{5.4cm} >{\raggedright\arraybackslash}p{2.6cm} >{\raggedright\arraybackslash}X}
\headrow \thead{Measure} & \thead{Value} & \thead{What changed it}\\
\midrule
Test-set accuracy before review & 67.3 percent & Initial clinical BERT pipeline \\
Test-set accuracy after review & 94.0 percent & Tabular model plus few-shot learning \\
Overall workflow quality & 87.5 percent & Evaluation standard, eight criteria \\
Efficiency score & 0.831 & Human, prompt, time and cost metrics \\
Regulatory analyst scores & 8.0, 9.1, 9.1 of 10 & Three manufacturers, recorded separately \\
Study duration, one author & 14 days & Dataset through FastAPI repository \\
\end{tabularx}
\end{apptable}
\vspace{-0.6cm}
\tabcap{Table 23. What the AI peer review changed, measured, in the glioblastoma\\study
whose method the pancreatic cancer work in this paper inherits.}

\figslot{graphviz-type / decision tree}{9.0}{Figure 24. One disagreement from
three model reports to five terminal\\dispositions, none of which is an average
and only one of which is silent.}

Figure 24 is the reason the disagreement in Table 23 is reported as three
separate scores rather than as a mean of 8.7. None of the five dispositions
averages the reports, and only one leaves a disagreement without an action.

\draftinstr{Stage 7: build Table 24 listing every instance of AI peer review
across the author's works, with columns Work, Date, Reviewers, What review
changed. Include the glioblastoma study, the code generation competition, the
NIH application v1 and v2, and the capitalization plan. Then close the section
by stating that AI peer reviewers are the only reviewers capable of handling
increasingly larger, more complex, and more frequent studies, drawing that claim
from the source study's abstract, introduction, limitations and conclusions.}

\begin{apptable}
\begin{tabularx}{\textwidth}{>{\raggedright\arraybackslash}p{5.0cm} >{\raggedright\arraybackslash}p{1.9cm} >{\raggedright\arraybackslash}p{2.6cm} >{\raggedright\arraybackslash}X}
\headrow \thead{Work} & \thead{Date} & \thead{Reviewers} & \thead{What the review changed}\\
\midrule
Glioblastoma patient matching ML & Nov 30, 2025 & 3 manufacturers & Model class, 67.3 to 94.0 percent \\
Code generation competition & Dec 22, 2025 & 16 models & Iterative learning on FDA adverse events \\
NIH application v1.0 & Jul 7, 2026 & 2 reviewers & Structure and mechanism fit \\
NIH application v2.0 & Jul 12, 2026 & 2 reviewers & Budget, milestones, review schedule \\
Capitalization plan & Aug 11, 2026 & 2 reviewers & Column widths, caption balance, figure centering \\
\end{tabularx}
\end{apptable}
\vspace{-0.6cm}
\tabcap{Table 24. Five instances of AI peer review across the author's deposited\\works,
and the specific change each review produced.}

\draftinstr{Stage 8: this is a main section. Match the character count of
sections 3, 4, 5 and 6, and verify every figure and table reference.}
